%\documentclass{book}\usepackage{sjtfonts}\usepackage{includegraphics}\usepackage{alltt}\usepackage{latexsym}
%\usepackage{array}
%\begin{document}\input{defs0}%		miradefs.tex 		    June 1994

% opening and closing a quoted Miranda section

\newcommand{\mira}{\noindent\hrulefill\nopagebreak\so}
\newcommand{\arim}{\nopagebreak\st\hrulefill\\}

% backslash in the Miranda environment
% Only feasible approach is to use \verb+\+
% in line!

%\newcommand{\bs}{$\mi\backslash\!\!$}
%\newcommand{\lbs}{\mbox{\verb+\+}}

% ignoring part of the input

\newcommand{\ignore}[1]{}

% a blank line in a Miranda section

\newcommand{\bl}{\mbox{\ }}

% Signalling a diffculty.

\definecolor{light-gray}{gray}{0.95}

\newcommand{\beware}[2]
 {\medskip\noindent\colorbox{light-gray}{\parbox{\textwidth}
 {\mbox{\large\textbf{#1}} \medskip\noindent\\ #2}\medskip\noindent}}

% Subscripting in Miranda text
\newcommand{\subscr}[2]{\mbox{${\mbox{\texttt{#1}}}_{\mbox{\texttt{#2}}}$}}

%\newcommand{\subscr}[2]{\mbox{${\mbox{\mi #1}}_{\mbox{\smtt #2}}$}}
\newcommand{\aone}{\subscr{a}{1}}
\newcommand{\atwo}{\subscr{a}{2}}
\newcommand{\ak}{\subscr{a}{k}}

\newcommand{\aoneone}{\subscr{a}{1,1}}

\newcommand{\eone}{\subscr{e}{1}}
\newcommand{\etwo}{\subscr{e}{2}}
\newcommand{\ei}{\subscr{e}{i}}
\newcommand{\er}{\subscr{e}{r}}

\newcommand{\fone}{\subscr{f}{1}}
\newcommand{\fl}{\subscr{f}{l}}

\newcommand{\gone}{\subscr{g}{1}}
\newcommand{\gtwo}{\subscr{g}{2}}
\newcommand{\gi}{\subscr{g}{i}}
\newcommand{\gr}{\subscr{g}{r}}

\newcommand{\hone}{\subscr{h}{1}}
\newcommand{\hl}{\subscr{h}{l}}

\newcommand{\pone}{\subscr{p}{1}}
\newcommand{\ptwo}{\subscr{p}{2}}
\newcommand{\pk}{\subscr{p}{k}}

\newcommand{\qone}{\subscr{q}{1}}
\newcommand{\qtwo}{\subscr{q}{2}}
\newcommand{\qk}{\subscr{q}{k}}

\newcommand{\rone}{\subscr{r}{1}}
\newcommand{\rtwo}{\subscr{r}{2}}

\newcommand{\tone}{\subscr{t}{1}}
\newcommand{\ttwo}{\subscr{t}{2}}
\newcommand{\tn}{\subscr{t}{n}}

\newcommand{\vone}{\subscr{v}{1}}
\newcommand{\vtwo}{\subscr{v}{2}}
\newcommand{\vn}{\subscr{v}{n}}

% for regular expressions

\newcommand{\stone}{\subscr{st}{1}}
\newcommand{\sttwo}{\subscr{st}{2}}
\newcommand{\stn}{\subscr{st}{n}}
\newcommand{\rthree}{\subscr{r}{3}}

\newcommand{\eps}{$\varepsilon$}
\newcommand{\trans}[1]{\stackrel{\mbox{\tt #1}}{\longrightarrow}}



% superscript in Miranda text

\newcommand{\superscr}[2]{\mbox{${\mbox{\mi #1}}^{\mbox{\smtt #2}}$}}

% Not equals in Miranda text

\newcommand{\noteq}{\makebox[0pt][l]{=}/}
\newcommand{\overslash}{\makebox[0pt][l]{/}}

% Definitions in complexity chapter

\newfont{\oldSym}{cmsy10}
\newcommand{\Th}{\boldmath$\oldSym\Theta$}
\newcommand{\pow}[1]{\^{}#1}
\newcommand{\leqq}{$\leq$}
\newcommand{\geqq}{$\geq$}
\newcommand{\lll}{$\ll$}
\newcommand{\eqq}{$\equiv$}
\newcommand{\ltwo}{\subscr{log}{2}}

% Indexing

\newcommand{\minx}[1]{\index{#1@{\ttfamily{}#1}}}
\setcounter{chapter}{17}\setcounter{page}{363}

\chapter{I/O programming}
\chapstart
\label{io-programming}
%\label{actions}
\index{I/O|(}


This chapter looks again at I/O in Haskell, particularly focussing on the \texttt{do} notation which we use to write I/O programs.  In the next chapter we'll see that we can write all sorts of different programs using \texttt{do} notation, including programs that manipulate state, that non-deterministically return more than one result, as well as programs that may fail, or raise exceptions. 
Underlying each of these is a \textbf{monad}. A monad provides the infrastructure for sequencing and naming used in the \texttt{do} notation, and we will show instances of the \texttt{Monad} class for lists, the \texttt{Maybe} type, a parsing monad and more.




\section{I/O programming}\index{interactive programs}

We first looked at I/O programming in Haskell using the \texttt{IO} types in Chapter \ref{io}. Something of type \texttt{IO t} is a \emph{program} which will perform I/O and return a value of type \texttt{t}. Among the I/O primitives of Haskell are
\begin{alltt}
putStr :: String -> IO ()
getLine :: IO String
\end{alltt}
The effect of running \texttt{putStr str} is to output the string \texttt{str} and then to terminate; it will return the value \texttt{()} which is (the only value) of type \texttt{()}.\minx{()} The effect of executing \texttt{getLine} is to get a line of input, and to return that line as its result.

Programs are built using the \texttt{do} notation to \emph{sequence statements} of type \texttt{IO t}. A program to read a line then output a message is given by
\begin{alltt}
readWrite :: IO ()

readWrite =
    do
      getLine
      putStrLn "one line read"
\end{alltt}
where the two statements are sequenced one after the other. It is also possible to \emph{name} the results of statements for use in the program, so that we can name the line read, and output it, like this:
\begin{alltt}
readEcho :: IO ()

readEcho =
    do
      line <-getLine
      putStrLn ("line read: " ++ line)
\end{alltt}
In this program \texttt{line <- getLine} names the result of the \texttt{getLine} and so it can be echoed in the next statement.

\subsection*{Adding a sequence of integers}

Now suppose we want to write an interactive program to sum integers
supplied one per line until zero is input. We will write an I/O program
\begin{alltt}\minx{sumInts}
sumInts :: Integer -> IO Integer
\end{alltt}
which is passed a starting value for the sum, and which returns the sum from there; so, to sum from scratch we call \texttt{sumInts 0}. This gives the program
\begin{alltt}
sumInts s
  = do n <- getInt
       if n==0 
          then return s
          else sumInts (s+n)
\end{alltt}
In writing the program there are two cases: we first read an integer value, using the \texttt{getInt} function defined in Chapter \ref{io}. If
we read zero then the result must be the sum so far, \texttt{s}; if not, we get the result by calling \texttt{sumInts} again, with the parameter set to \texttt{(s+n)}.



The program is written in the style we saw in Chapter \ref{io}: the function is \emph{tail-recursive} -- being called as the last statement in the program -- and the \emph{loop data} is used to carry the ``sum so far'': here \texttt{s} and \texttt{s+n}. 



It is interesting with compare the definition of \texttt{sumInts} with the recursion in
\begin{alltt}\minx{sumAcc}
sumAcc s [] = s
sumAcc s (n:ns) 
  = if n==0
       then s
       else sumAcc (s+n) ns
\end{alltt}
Here we use the first variable to \texttt{sumAcc} as an \emph{accumulator} where the ``result so far'' can be held. 

We can also put the \texttt{sumInts} program inside a `wrapper' which explains its
purpose and prints the sum at the end.
\begin{alltt}\minx{sumInteract}
sumInteract :: IO ()
sumInteract
  = do putStrLn "Enter integers one per line"
       putStrLn "These will be summed until zero is entered"
       sum <- sumInts 0
       putStr "The sum is "
       print sum
\end{alltt}


\begin{exercises}

\item
Compare the definition of \texttt{sumInts} above with this definition:
\begin{alltt}
sumInts :: IO Integer
sumInts
  = do n <- getInt
       if n==0 
          then return 0
          else (do m <- sumInts
                   return (n+m))
\end{alltt}
Which will have the better performance, do you think? Wny?

\item
Give a definition of the function
\begin{alltt}
fmap :: (a -> b) -> IO a -> IO b
\end{alltt}
the effect of which is to transform an interaction by applying the function to its
result. You should define it using the \texttt{do} construct.

\item
Define the function
\begin{alltt}
repeat ::  IO Bool -> IO () -> IO ()
\end{alltt}
so that \texttt{repeat test oper} has the effect of repeating \texttt{oper} until
the condition \texttt{test} is \texttt{True}.

\item
Define the higher-order function \texttt{while} in which the condition and the operation
work over values of type \texttt{a}. Its type should be
\begin{alltt}
whileG :: (a -> IO Bool) -> (a -> IO a) -> (a -> IO a)
\end{alltt}

\item 
Using the function \texttt{whileG} or otherwise, define an interaction which
reads a number, \texttt{n} say, and then reads a further \texttt{n} numbers and
finally returns their average.

\item
Modify your answer to the previous question so that if the end of file is
reached before \texttt{n} numbers have been read, a message to that effect is
printed.

\item
Define a function
\begin{alltt}
accumulate :: [IO a] -> IO [a]
\end{alltt}
which performs a sequence of interactions and accumulates their result in a
list. Also give a definition of the function
\begin{alltt}
sequence :: [IO a] -> IO ()
\end{alltt}
which performs the interactions in turn, but discards their results. Finally,
show how you would sequence a series, passing values from one to the next:
\begin{alltt}
seqList :: [a -> IO a] -> a -> IO a
\end{alltt}
What will be the result on an empty list?
\end{exercises}

\section{Further I/O}
\label{furtherIO}

In this section we survey further features of Haskell I/O, defined in the \texttt{System.IO} module.

\subsection*{Interaction at the terminal}\index{interactive programs}

We have seen that we can read from the terminal -- the `standard input' -- and write to
the screen -- the `standard output'. Terminal input can be configured to work in different ways, depending on the way that the input is \texttt{buffered}. 

The default is that input is unbuffered, and so each character is available to the I/O program immediately it is typed; a disadvantage of this mode is that it is impossible to use \texttt{Ctrl-D} to signal ``end of file'' in the interactive input. Setting to line buffering, where input is assembled into lines before being available to the I/O program, remedies this problem. To set the buffering modes we use \texttt{hSetBuffering} as in this example:
\begin{alltt}\minx{copyInteract}
copyInteract :: IO ()

copyInteract = 
    do
      hSetBuffering stdin LineBuffering
      copyEOF
      hSetBuffering stdin NoBuffering
\end{alltt}
where we can see that the buffering mode can be changed within an I/O program. The \texttt{copyEOF} program itself is given by 
\begin{alltt}\minx{copyEOF}
copyEOF :: IO ()

copyEOF = 
    do 
      eof <- isEOF
      if eof  
        then return () 
        else do line <- getLine 
                putStrLn line
                copyEOF
\end{alltt}
where \texttt{isEOF ::\ IO Bool} is an I/O program to determine whether or not the end of standard input has been reached. 

Give this a try within GHCi: you will find that you can use \texttt{Ctrl-D} to terminate your input if you run \texttt{copyInteract}; if on the other hand you run \texttt{copyEOF} directly, you will need to interrupt the program using \texttt{Ctrl-C}.

\subsection*{File I/O}
\index{I/O!files}
As well as reading and writing to a terminal, the Haskell I/O model also provides for reading from and
writing and appending to files, by means of the functions
\begin{alltt}\minx{readFile}\minx{writeFile}\minx{appendFile}
readFile   :: FilePath -> IO String
writeFile  :: FilePath -> String -> IO ()
appendFile :: FilePath -> String -> IO ()
\end{alltt} 
where
\begin{alltt}\minx{FilePath}
type FilePath = String
\end{alltt}                                                                               
and files are specified by the text strings appropriate to the
implementation in question.

\subsection*{Errors}
\index{I/O!errors}

I/O programs can raise errors, which belong to the system-dependent
\texttt{data} type \texttt{IOError}. The function
\begin{alltt}\minx{ioError}
ioError :: IOError -> IO a
\end{alltt}
builds an I/O action which fails giving the appropriate error, and the program
\begin{alltt}\minx{catch}
catch :: IO a -> (IOError -> IO a) -> IO a
\end{alltt}
will catch an error raised by the first argument and
handle it using the second argument, which gives a handler -- that is an
action of type \texttt{IO a} -- for each possible \texttt{IOError}.

\subsection*{Input and output as lazy lists}
\index{I/O!as lazy lists}

An alternative view of I/O programs, popular in earlier lazy functional 
programming languages, was to see the input and output as
\texttt{Strings}, that is as lists of characters. Under that model an
I/O program is a function
\begin{alltt}
listIOprog :: String -> String
\end{alltt}
This obviously makes sense in a `batch' program, where all the input is
read before any output is produced, but in fact it also works for
interactive programs where input and output are interleaved, if the
language is \textbf{lazy}. This is because in a lazy language we can
begin to print the result of a computation -- the output of the
interactive program here -- before the argument -- the interactive input --
is fully evaluated. As an example, repeatedly to reverse lines of input
under this model one can write
\begin{alltt}
listIOprog  = unlines . map reverse . lines
\end{alltt}
The drawback of this approach is in scaling it up. It is often difficult
to predict in advance the way in which the input and output are
interleaved: often output comes after it is expected, and sometimes even
before; the \texttt{IO} approach in Haskell avoids such problems.
Nevertheless, support for this style is available, using
\begin{alltt}\minx{getContents}
getContents :: IO String
\end{alltt}
a primitive to get the contents of the standard input, and which is used
in this function, also defined in \texttt{System.IO}\minx{System.IO}
\begin{alltt}\minx{interact}
interact :: (String -> String) -> IO ()
interact f = do s <- getContents
                putStr (f s)
\end{alltt}
Try out the program \texttt{interact listIOprog} in GHCi to see for yourself that the input and output lines are interleaved one-by-one just as you would expect.

\subsection*{Generating random numbers}\index{random numbers}

We introduced the \texttt{randomStrategy} in Chapter \ref{io}, but pointed out there that it should properly be programmed using a monad, since it is a function which returns different values on different calls. We generate a random integer less than \texttt{n} based on the current time, which is accessible within the \texttt{IO} monad:
\begin{alltt}\minx{randomInt}
randomInt :: Integer -> IO Integer
randomInt n = 
    do
      time <- getCurrentTime
      return ( (`rem` n) $ read $ take 6 $ 
                         formatTime defaultTimeLocale "%q" time)
      
randInt :: Integer -> Integer
randInt = unsafePerformIO . randomInt 
      
sRandom :: Strategy
sRandom _ = convertToMove (randInt 3)
\end{alltt}
Here we use the \texttt{unsafePerformIO} function to extract the result from the \texttt{IO} monad. As its name suggests this is (potentially) unsafe, and should be used with care; some would go further, and say that it should never be used.

If we want to avoid \texttt{unsafePerformIO} we need to look again at the design of our case study. In Chapter \ref{io} defined strategies to belong to the type 
\begin{alltt}\minx{Strategy}
type Strategy = [Move] -> Move
\end{alltt}
if we want to include monadic functions, then we should redefine them to be monadic, like this
\begin{alltt}
type StrategyM = [Move] -> IO Move
\end{alltt}
We leave the reimplementation of the case study to use monadic strategies as an exercise for the reader.

\subsection*{The \texttt{System.IO} library}
\minx{System.IO}

There is much more to the  \texttt{System.IO} library than we are able to cover here; see the system documentation for more details.

\begin{exercises}

\item
Write a file-handling version of the program
\texttt{sumInts}, which reads the integer sequence from a file.

\item
Write a lazy-list version of the program \texttt{sumInts}, which you can then run using \texttt{interact}.

\item{[Harder]}
Write a lazy-list version of the calculator program.

\item{[Harder]}
Reimplement the Rock - Paper - Scissors case study to use the monadic strategies in \texttt{StrategyM}.

\end{exercises}

\section{The calculator}
\label{interactiveCalc}
\index{calculator|(}

The ingredients of the calculator are contained in three places in the text.
\begin{titemize}
\item In Section \ref{recTypes} we saw the introduction of the algebraic type of
expressions, \texttt{Expr}, which we subsequently revised in Section
\ref{parsing}, giving
\begin{alltt}
data Expr = Lit Integer | Var Var | Op Ops Expr Expr\minx{Expr}
data Ops  = Add | Sub | Mul | Div | Mod\minx{Ops}
type Var  = Char
\end{alltt}
We revise the evaluation of expressions after discussing the store below. 
\item In Chapter \ref{adt} we introduced the abstract type
\texttt{Store}\minx{Store}, which we use to model the values of the variables
currently held. The signature of the abstract data type is
\begin{alltt}
initial :: Store 
value   :: Store -> Var -> Integer
update  :: Store -> Var -> Integer -> Store
\end{alltt}
\item In Section \ref{parsing} we looked at how to parse expressions and
commands, 
\begin{alltt}\minx{Command}
data Command = Eval Expr | Assign Var Expr | Null
\end{alltt}
and defined the ingredients of the function
\begin{alltt}
commLine :: String -> Command\minx{commLine}
\end{alltt}
which is used to parse each line of input into a \texttt{Command}. For instance,
\begin{alltt}
commLine "(3+x)"   = (Eval (Op Add (Lit 3) (Var 'x')))
commLine "x:(3+x)" = (Assign 'x' (Op Add (Lit 3) (Var 'x')))
commLine "something unparsable" = Null
\end{alltt} 
\end{titemize}
Expressions are evaluated by
\begin{alltt}
eval :: Expr -> Store -> Integer\minx{eval}

eval (Lit n) st = n
eval (Var v) st = value st v
eval (Op op e1 e2) st
  = opValue op v1 v2
    where
    v1 = eval e1 st
    v2 = eval e2 st
\end{alltt}
where the \texttt{opValue} function  of type \texttt{Ops->Integer->Integer->Integer}
interprets each operator, such as \texttt{Add}, as the corresponding function,
like \texttt{(+)}. 

What is the effect of a command? An expression should return the value
of the expression in the current store; an assignment will change the
store, and a null command will do nothing. We therefore define a
function which returns both a value and a store,
\begin{alltt}
command :: Command -> Store -> (Integer,Store)\minx{command}

command Null st     = (0 , st)
command (Eval e) st = (eval e st , st)
command (Assign v e) st 
  = (val , newSt)
    where
    val   = eval e st
    newSt = update st v val
\end{alltt}
A single step of the calculator will take a starting \texttt{Store}, and
read an input line, evaluate the command in the line, print some output
and finally return an updated \texttt{Store},
\begin{alltt}\minx{calcStep}
calcStep :: Store -> IO Store

calcStep st
  = do line <- getLine
       let comm        = commLine line\hfill(1)
       let (val,newSt) = command comm st\hfill(2)
       print val
       return newSt
\end{alltt}
In lines \texttt{(1)} and \texttt{(2)}
of the definition of \texttt{calcStep} we see an example of the use
of \texttt{let} within a \texttt{do} expression. Line
\texttt{(1)},
\begin{alltt}\minx{let}
let comm        = commLine line
\end{alltt}
gives \texttt{comm} the value of parsing the \texttt{line}, and this is
subsequently used in \texttt{(2)},
\begin{alltt}
let (val,newSt) = command comm st
\end{alltt}
which simultaneously gives \texttt{val} and \texttt{newSt} the value of
the expression read and the new state. In the lines that
follow the \texttt{let}s, the value \texttt{val} is printed and the new state \texttt{newSt} is returned as the
overall result of the interaction. 
A sequence of calculator steps is given by
\begin{alltt}\minx{calcSteps}
calcSteps :: Store -> IO ()

calcSteps st =
    do
      eof <- isEOF
      if eof
         then return ()
         else do newSt <- calcStep st
                 calcSteps newSt
\end{alltt}
The main I/O program for the calculator is given by starting off
\texttt{calcSteps} with the \texttt{initial} store, operating in modified buffering mode
\begin{alltt}\minx{mainCalc}
mainCalc :: IO ()
mainCalc = 
    do
      hSetBuffering stdin LineBuffering
      calcSteps initial
      hSetBuffering stdin NoBuffering
\end{alltt}
In the exercises various extensions and modifications of the calculator
program are discussed.


\begin{exercises}

\item
How would you add initial and final messages to the output of the calculator?

\item
Discuss how you would have to modify the system to allow variables to have
arbitrarily long names, consisting of letters and numbers, starting with a
letter.

\item
How would you extend the calculator to deal with decimal floating-point numbers
as well as
integers?

\item
Discuss how you would modify the calculator so that it could read input
commands split over more than one line. You will need to decide how this
sort of split is signalled by the user -- maybe by \texttt{\bs{}} at the end of the line --
and how to modify the interaction program to accommodate this.
Alternatively, you might let the user do this without signalling; can
you modify the program to do that?

\item
How would you modify the parser so that `white space' is permitted in
the input commands, as in the example
\begin{alltt}
"    x : (2\bs{t}+3)     "
\end{alltt}
which parses to the \texttt{Command}
\begin{alltt}
(Assign 'x' (Op Add (Lit 2) (Lit 3)))
\end{alltt}
\index{calculator|)}
\end{exercises}
\index{I/O|)}



\begin{summary}%

In this chapter we have seen more examples of how the \texttt{do} notation can be used for I/O programming. The mechanism which underlies the \texttt{do} notation is the \texttt{Monad} class, and we return to that in the next chapter.

\end{summary}

%\end{document}
